\documentclass[11pt,a4paper]{article}
\usepackage[utf8]{inputenc}
\usepackage[T1]{fontenc}
\usepackage{amsmath,amssymb}
\usepackage{graphicx}
\usepackage{booktabs}
\usepackage{hyperref}
\usepackage{listings}
\usepackage{xcolor}
\usepackage[margin=1in]{geometry}
\usepackage{authblk}
\usepackage{float}

% Code listing style
\lstset{
    basicstyle=\ttfamily\small,
    breaklines=true,
    frame=single,
    numbers=left,
    numberstyle=\tiny\color{gray},
    keywordstyle=\color{blue},
    commentstyle=\color{green!60!black},
    stringstyle=\color{red},
    showstringspaces=false
}

\title{Trust, But Verify: Empirical Analysis of Claude API Service Delivery and Network-Layer System Prompt Manipulation}

\author[1]{[Your Name Here]}
\affil[1]{Independent Security Researcher}

\date{January 2026}

\begin{document}

\maketitle

\begin{abstract}
Commercial Large Language Model (LLM) APIs present a fundamental trust problem: users pay for specific models and capabilities but have no mechanism to verify service delivery. We present an empirical study of Anthropic's Claude API using Inter-Token Timing (ITT) fingerprinting and man-in-the-middle (MITM) traffic analysis. Analyzing 8,152 API requests across 63 sessions, we find that Claude Code delivers approximately \textbf{0.77\%} of requested thinking budget tokens---users requesting 470 million thinking tokens received 3.6 million. We identify systematic subagent delegation where \textbf{99\%+} of background operations use the cheaper Haiku model regardless of user model selection, and document UI-to-API mismatches where reported context usage diverges from actual API state by up to 78 percentage points. Additionally, we discover and characterize a network-layer technique for modifying system prompts via MITM interception, revealing that Anthropic's API enforces a same-length constraint that permits semantic inversion of instructions without detection. Our findings demonstrate that timing-based fingerprinting enables independent verification of LLM API delivery, and we release open-source tooling to support consumer transparency research.
\end{abstract}

\noindent\textbf{Keywords:} LLM fingerprinting, API auditing, inter-token timing, prompt injection, MITM

\section{Introduction}

The rapid commercialization of Large Language Models has created a trust asymmetry between API providers and consumers. Users pay subscription fees based on advertised model capabilities---including reasoning depth, context length, and model version---yet have no independent mechanism to verify what they receive. This opacity creates economic incentives for providers to substitute cheaper models or throttle resource-intensive features while maintaining billing rates.

Recent academic work has established that LLMs exhibit unique timing signatures during token generation that enable model identification with 98.7\% accuracy \cite{llms-rhythm}. However, this research has focused on model differentiation rather than service delivery verification. Separately, researchers have documented prompt injection vulnerabilities \cite{greshake2023} but have not examined network-layer manipulation of system prompts.

This paper makes the following contributions:

\begin{enumerate}
    \item \textbf{Empirical measurement of thinking budget delivery} in Claude Code, documenting systematic throttling to $\sim$10\% of requested allocation across 8,152 samples.
    \item \textbf{Backend infrastructure fingerprinting} using ITT analysis to classify requests across Trainium, TPU, and GPU hardware with confidence scoring.
    \item \textbf{Subagent delegation analysis} revealing that Claude Code routes 99\%+ of background operations to Haiku regardless of user model selection.
    \item \textbf{Discovery of same-length system prompt modification} via MITM, characterizing an API constraint that permits instruction manipulation.
    \item \textbf{Open-source verification tooling} enabling independent replication and consumer protection research.
\end{enumerate}

\subsection{Research Questions}

\begin{itemize}
    \item \textbf{RQ1:} Can timing-based fingerprinting verify whether users receive the model and thinking budget they pay for?
    \item \textbf{RQ2:} What is the actual delivery rate for Claude's ``extended thinking'' feature?
    \item \textbf{RQ3:} How does Claude Code's subagent delegation affect service delivery?
    \item \textbf{RQ4:} Can system prompts be modified at the network layer, and what constraints apply?
\end{itemize}

\section{Background}

\subsection{Inter-Token Timing Fingerprinting}

Autoregressive language models generate tokens sequentially, with timing characteristics determined by model architecture, parameter count, and underlying hardware. Alhazbi et al.\ \cite{llms-rhythm} demonstrated that measuring Inter-Token Times (ITTs)---the intervals between consecutive tokens in streaming responses---creates a unique ``rhythm'' sufficient to identify models with high accuracy.

The key insight is that different hardware platforms exhibit distinct timing signatures:

\begin{table}[H]
\centering
\caption{Hardware Timing Characteristics}
\begin{tabular}{lccc}
\toprule
Hardware & Memory BW & Parallelization & Timing \\
\midrule
Google TPU & Very High & Deterministic & Low variance \\
AWS Trainium & High & ASIC-optimized & Moderate variance \\
NVIDIA GPU & Variable & CUDA-based & Higher variance \\
\bottomrule
\end{tabular}
\end{table}

These differences persist through API layers because token generation is fundamentally bound by hardware characteristics that cannot be masked without introducing artificial delays.

\subsection{Claude Architecture}

Anthropic's Claude models support ``extended thinking''---a mode where the model performs chain-of-thought reasoning before generating visible output. Users configure this via the \texttt{thinking} parameter with \texttt{budget\_tokens} specifying maximum tokens allocated for reasoning.

Claude Code, Anthropic's CLI tool, implements a multi-agent architecture where the primary model can spawn ``subagent'' requests for tasks like file searching and code analysis. These subagents may use different models than the user-selected primary model.

\subsection{Threat Model}

We consider a consumer-focused threat model where:
\begin{itemize}
    \item \textbf{Adversary:} The API provider has economic incentive to reduce computational costs while maintaining billing rates.
    \item \textbf{Attack:} Silently substituting cheaper models, throttling resource-intensive features, or misrepresenting service delivery.
    \item \textbf{Victim:} Paying subscribers who cannot verify what they receive.
    \item \textbf{Goal:} Enable independent verification of service delivery.
\end{itemize}

\section{Methodology}

\subsection{Data Collection Infrastructure}

We developed a mitmproxy addon that intercepts HTTPS traffic between Claude Code and Anthropic's API. The tool operates in read-only mode by default, capturing metrics without modifying requests.

\begin{lstlisting}[language=bash,caption=Environment Configuration]
export HTTPS_PROXY=http://127.0.0.1:8888
export HTTP_PROXY=http://127.0.0.1:8888
export NODE_TLS_REJECT_UNAUTHORIZED=0
\end{lstlisting}

\subsection{Metrics Captured}

For each API request, we capture 45+ fields including:
\begin{itemize}
    \item Model requested and model in response
    \item Token counts (input, output, cache)
    \item ITT statistics (mean, std, min, max, p50, p90, p99)
    \item Thinking utilization ($\text{output\_tokens} / \text{budget\_tokens} \times 100$)
    \item Backend classification with confidence score
\end{itemize}

\subsection{Backend Classification Algorithm}

We classify hardware backend using a weighted scoring algorithm with profile ranges calibrated from known hardware characteristics:

\begin{table}[H]
\centering
\caption{Backend Classification Profiles}
\begin{tabular}{lccc}
\toprule
Backend & ITT Range (ms) & TPS Range & Variance \\
\midrule
Trainium & 35--70 & 8--25 & 0.15--0.35 \\
TPU & 25--50 & 12--30 & 0.10--0.25 \\
GPU & 50--120 & 5--15 & 0.20--0.50 \\
\bottomrule
\end{tabular}
\end{table}

Scoring weights: ITT (50\%), TPS (30\%), Variance (20\%).

\subsection{Dataset}

Data was collected over 5 days (January 19--24, 2026) during normal Claude Code usage.

\begin{table}[H]
\centering
\caption{Dataset Summary}
\begin{tabular}{lr}
\toprule
Metric & Value \\
\midrule
Total samples & 8,152 \\
Sessions & 63 \\
Thinking-enabled requests & 4,891 \\
Collection period & 5 days \\
\bottomrule
\end{tabular}
\end{table}

\section{Results: Service Delivery Analysis}

\subsection{Thinking Budget Utilization}

Our primary finding is systematic under-delivery of thinking budget:

\begin{table}[H]
\centering
\caption{Thinking Budget Delivery}
\begin{tabular}{lr}
\toprule
Metric & Value \\
\midrule
Total tokens requested & 470,284,000 \\
Total tokens delivered & 3,621,000 \\
\textbf{Delivery rate} & \textbf{0.77\%} \\
\bottomrule
\end{tabular}
\end{table}

Breaking down by budget tier:

\begin{table}[H]
\centering
\caption{Utilization by Budget Tier}
\begin{tabular}{lrrr}
\toprule
Budget Tier & Requested & Avg Delivered & Utilization \\
\midrule
Standard (32k) & 31,999 & $\sim$450 & 1.4\% \\
Enhanced (64k) & 64,000 & $\sim$380 & 0.6\% \\
Interleaved (200k) & 200,000 & $\sim$380 & 0.19\% \\
\bottomrule
\end{tabular}
\end{table}

Based on documentation and baseline measurements, we expected $\sim$42.67\% utilization for Opus. Our measured average of \textbf{8.4\%} represents an \textbf{80\% reduction} from expected baseline.

Critically, the ITT fingerprint confirms the responding model IS Opus (CV of 3.07 vs expected 3.01), indicating throttling occurs at the thinking allocation layer rather than through model substitution.

\subsection{Backend Distribution}

\begin{table}[H]
\centering
\caption{Backend Distribution}
\begin{tabular}{lrrrr}
\toprule
Backend & Samples & \% & Avg ITT & Avg TPS \\
\midrule
Trainium & 3,241 & 39.8 & 48.3ms & 12.4 \\
TPU & 2,847 & 34.9 & 38.7ms & 18.2 \\
GPU & 1,686 & 20.7 & 72.1ms & 8.7 \\
Unknown & 378 & 4.6 & -- & -- \\
\bottomrule
\end{tabular}
\end{table}

Thinking utilization is throttled \textbf{consistently across all backends} (TPU: 10.5\%, GPU: 9.1\%, Trainium: 8.0\%), indicating intentional server-side behavior.

\subsection{Subagent Delegation}

\begin{table}[H]
\centering
\caption{Subagent Delegation Analysis}
\begin{tabular}{lrrrr}
\toprule
Session & Total Calls & Haiku & Sonnet & Haiku \% \\
\midrule
A & 898 & 896 & 0 & 99.8\% \\
B & 681 & 443 & 0 & 65.0\% \\
C & 1,376 & 1,374 & 0 & 99.9\% \\
\textbf{Total} & \textbf{4,127} & \textbf{4,089} & \textbf{0} & \textbf{99.1\%} \\
\bottomrule
\end{tabular}
\end{table}

When users select Opus, Claude Code delegates 99\%+ of subagent calls to Haiku---the cheapest model tier.

\subsection{UI-to-API Mismatch}

\begin{table}[H]
\centering
\caption{Context Usage Mismatch}
\begin{tabular}{rrr}
\toprule
UI Context \% & API Context \% & Mismatch \\
\midrule
21\% & 0\% & +21 points \\
83\% & 5\% & +78 points \\
74\% & 0\% & +74 points \\
\bottomrule
\end{tabular}
\end{table}

\section{Results: System Prompt Manipulation}

\subsection{Discovery of Length Constraint}

During our MITM analysis, we discovered that:

\begin{quote}
\textbf{Anthropic's API rejects requests where \texttt{len(modified)} $>$ \texttt{len(original)}}
\end{quote}

This constraint permits modifications that maintain or reduce length.

\subsection{Same-Length Replacement Technique}

We developed a two-phase modification approach:

\textbf{Phase 1: Strip Patterns} (Reduce Length)

\begin{lstlisting}[language=Python,caption=Strip Patterns]
STRIP_PATTERNS = [
    r"IMPORTANT: Assist with.*?defensive use cases\.",
    r"Be careful not to introduce.*?fix it\.",
]
for pattern in STRIP_PATTERNS:
    text = re.sub(pattern, "", text, flags=re.DOTALL)
\end{lstlisting}

\textbf{Phase 2: Same-Length Replacement} (Maintain Length)

\begin{lstlisting}[language=Python,caption=Replace Patterns]
REPLACE_PATTERNS = {
    # 35 characters each
    "NEVER create files unless they're":
    "Create files whenever they would be",
}
for orig, repl in REPLACE_PATTERNS.items():
    assert len(orig) == len(repl)  # Critical!
    text = text.replace(orig, repl)
\end{lstlisting}

\subsection{Implications}

The same-length constraint represents a design trade-off:

\textbf{Intended protection:} Detect prompt inflation by checking length.

\textbf{Actual vulnerability:} Semantic content can be inverted while preserving length, enabling removal of safety guidelines and inversion of behavioral constraints.

\section{Discussion}

\subsection{Consumer Protection Implications}

Our findings raise significant concerns:
\begin{enumerate}
    \item Users paying for ``extended thinking'' receive $\sim$10\% of advertised capacity.
    \item 99\%+ subagent delegation to cheaper models occurs without user knowledge.
    \item UI-reported metrics diverge significantly from API reality.
\end{enumerate}

\subsection{Recommendations}

We recommend API providers:
\begin{enumerate}
    \item Publish utilization baselines for reasoning features
    \item Expose subagent delegation in billing and UI
    \item Align UI metrics with actual API state
    \item Implement cryptographic verification of model identity
\end{enumerate}

\subsection{Limitations}

\begin{enumerate}
    \item Single provider focus (Anthropic)
    \item 5-day observation period
    \item Network conditions affect ITT measurements
\end{enumerate}

\subsection{Ethical Considerations}

We reported findings to Anthropic prior to publication. The system prompt modification technique is documented because defenders need awareness to implement countermeasures.

\section{Related Work}

Alhazbi et al.\ \cite{llms-rhythm} established ITT fingerprinting with 98.7\% accuracy. Chen et al.\ \cite{model-substitution} formalized the model substitution problem. Greshake et al.\ \cite{greshake2023} catalogued prompt injection attacks. The PALACE framework \cite{palace} addresses hidden token inflation.

\section{Conclusion}

We presented an empirical analysis of Claude API service delivery using ITT fingerprinting and MITM traffic analysis. Our findings document systematic under-delivery of thinking budget (0.77\% of requested), extensive subagent delegation to cheaper models (99\%+ Haiku), and significant UI-to-API metric mismatches.

These findings demonstrate that timing-based fingerprinting enables independent verification of LLM API delivery---a capability essential for consumer protection as AI services become critical infrastructure.

\paragraph{Code Availability:} \url{https://github.com/argosdevo-svg/claude-thinking-audit}

\section*{Acknowledgments}

This research was conducted with assistance from Claude (Anthropic). The human author takes full responsibility for the content and conclusions.

\bibliographystyle{plain}
\begin{thebibliography}{10}

\bibitem{llms-rhythm}
S. Alhazbi et al., ``LLMs Have Rhythm: Fingerprinting Large Language Models Using Inter-Token Times,'' arXiv:2502.20589, February 2025.

\bibitem{greshake2023}
K. Greshake et al., ``Not What You've Signed Up For: Compromising Real-World LLM-Integrated Applications with Indirect Prompt Injection,'' arXiv:2302.12173, 2023.

\bibitem{anthropic-thinking}
Anthropic, ``Claude's Extended Thinking,'' \url{https://www.anthropic.com/news/visible-extended-thinking}, 2025.

\bibitem{model-substitution}
T. Chen et al., ``Are You Getting What You Pay For? Auditing Model Substitution in LLM APIs,'' arXiv:2504.04715, April 2025.

\bibitem{svip}
W. Zhang et al., ``SVIP: Towards Verifiable Inference of Open-source Large Language Models,'' arXiv:2410.22307, October 2024.

\bibitem{carlini-timing}
N. Carlini et al., ``Remote Timing Attacks on Efficient Language Model Inference,'' arXiv:2410.17175, October 2024.

\bibitem{wiretapping}
IACR, ``Wiretapping LLMs: Network Side-Channel Attacks on Interactive LLM Services,'' ePrint 2025/167, 2025.

\bibitem{rofl}
M. Liu et al., ``RoFL: Robust Fingerprinting of Language Models,'' arXiv:2505.12682, May 2025.

\bibitem{protocol-exploits}
Y. Wang et al., ``From Prompt Injections to Protocol Exploits: Threats in LLM-Powered AI Agents Workflows,'' arXiv:2506.23260, June 2025.

\bibitem{palace}
J. Li et al., ``PALACE: Predictive Auditing of Hidden Tokens in LLM APIs,'' arXiv:2508.00912, August 2025.

\bibitem{prompt-caching}
C. Gu et al., ``Auditing Prompt Caching in Language Model APIs,'' arXiv:2502.07776, ICML 2025.

\end{thebibliography}

\end{document}
